% =====================================================================
\chapter{Σχεδίαση του Συστήματος}
\label{ch:design}
% =====================================================================

\section{Απαιτήσεις}
\label{sec:requirements}

\subsection{Λειτουργικές απαιτήσεις}

\begin{requirement}[Κλειστή πύλη]\label{req:closed}
Κατά την εκκίνηση, κάθε ριζικός διανομέας λαμβάνει \code{authorized\_default=0}.
Οι συσκευές που ήταν ήδη συνδεδεμένες παραμένουν ανεπηρέαστες.
\end{requirement}

\begin{requirement}[Ανάλυση πριν από την ερώτηση]\label{req:analysis}
Για κάθε νέα συσκευή παράγεται αναφορά με ευρήματα σοβαρότητας
\sevC, \sevW{} ή \sevN, πριν ερωτηθεί ο χρήστης.
\end{requirement}

\begin{requirement}[Ανθρώπινη απόφαση]\label{req:human}
Μια συσκευή ενεργοποιείται οριστικά μόνο με απάντηση του χρήστη (τερματικό ή
γραφικός διάλογος), με αποθηκευμένη εμπιστοσύνη «πάντα», ή με ρητή εξαίρεση
ασφαλείας (εσωτερική θύρα, \code{--allow-port}).
\end{requirement}

\begin{requirement}[Μνήμη]\label{req:memory}
Το σύστημα θυμάται κάθε συσκευή που έχει δει, τις αποφάσεις για αυτήν και ένα
αποτύπωμα των περιγραφέων της.
\end{requirement}

\begin{requirement}[Αδυναμία απάντησης]\label{req:no-answer}
Όταν δεν υπάρχει διαθέσιμος άνθρωπος (κλειδωμένη οθόνη, οθόνη σύνδεσης), οι νέες
συσκευές \emph{κρατούνται} απενεργοποιημένες και ερωτώνται αργότερα. Η απουσία
απάντησης δεν μετατρέπεται ποτέ σε αποδοχή.
\end{requirement}

\subsection{Μη λειτουργικές απαιτήσεις}

\begin{requirement}[Αποφυγή κλειδώματος]\label{req:lockout}
Η αρχική κατάσταση των διανομέων αποκαθίσταται σε κάθε έξοδο, κατάρρευση ή
σήμα. Υπάρχουν μηχανισμοί διάσωσης (\en{watchdog}, αρχείο πανικού,
\code{--release}).
\end{requirement}

\begin{requirement}[Ελάχιστο προνόμιο]\label{req:privsep}
Η ανάλυση μη έμπιστων δεδομένων (περιγραφείς, συμβολοσειρές, μέσα) εκτελείται σε
διεργασία χωρίς προνόμια.
\end{requirement}

\begin{requirement}[Φορητότητα]\label{req:portable}
Λειτουργία σε αμετάβλητο πυρήνα \en{Linux} με \en{Python}~3.10+ και μόνη
εξωτερική εξάρτηση το \en{pyudev}.
\end{requirement}

\begin{requirement}[Ελεγξιμότητα]\label{req:testable}
Η σουίτα ελέγχων εκτελείται χωρίς \en{root}, χωρίς υλικό \en{USB} και χωρίς
\en{pyudev}.
\end{requirement}

\section{Επισκόπηση της αρχιτεκτονικής}
\label{sec:architecture}

Το Σχήμα~\ref{fig:architecture} παρουσιάζει την αρχιτεκτονική του συστήματος
στη λειτουργία διαχωρισμού προνομίων (\code{--privsep}). Το σύστημα αποτελείται
από τρεις διεργασίες:

\begin{enumerate}
  \item \textbf{Η πύλη} (\en{gate}, \en{root}): μικρή διεργασία που γράφει τα
        αρχεία \code{authorized}, ανοίγει κόμβους συσκευών και περνά τους
        περιγραφείς αρχείων στον αναλυτή.
  \item \textbf{Ο αναλυτής} (\en{analyzer}, χρήστης \code{probolos}): εκτελεί
        τους κανόνες, την καραντίνα, το ιστορικό και τη ροή αποφάσεων.
  \item \textbf{Ο πράκτορας επιφάνειας εργασίας} (\en{agent}, χρήστης
        επιφάνειας εργασίας): εμφανίζει τους διαλόγους και επιστρέφει την
        απάντηση.
\end{enumerate}

\begin{figure}[htbp]
\centering
\begin{tikzpicture}[
  proc/.style={draw, thick, rounded corners=3pt, minimum width=2.9cm,
               minimum height=1.3cm, align=center, font=\small},
  kern/.style={draw, fill=gray!12, minimum width=14cm, minimum height=0.9cm,
               font=\small, align=center},
  lab/.style={font=\scriptsize, align=center},
  >={Stealth[length=2.4mm]}]

  \node[proc, fill=critical!10] (gate) at (0,0) {\textbf{Πύλη}\\\footnotesize(\en{root})\\\footnotesize\code{gate\_server.py}};
  \node[proc, fill=notice!10]   (an)   at (5.6,0) {\textbf{Αναλυτής}\\\footnotesize(\code{probolos})\\\footnotesize\code{daemon.py}};
  \node[proc, fill=okgreen!10]  (ag)   at (11.2,0) {\textbf{Πράκτορας}\\\footnotesize(χρήστης)\\\footnotesize\code{agent.py}};

  \node[kern] (k) at (5.6,-2.6) {Πυρήνας \en{Linux}: \code{/sys/bus/usb}\ \textbullet\ \code{/dev/input/eventN}\ \textbullet\ \code{/dev/sdX}\ \textbullet\ \en{uevents}};

  \draw[<->, thick] (gate) -- node[above, lab]{\en{SEQPACKET}\\\en{socketpair}} node[below, lab]{\en{SCM\_RIGHTS}} (an);
  \draw[<->, thick] (an) -- node[above, lab]{\en{AF\_UNIX}} node[below, lab]{\en{SO\_PEERCRED}} (ag);
  \draw[->] (gate.south) -- node[left, lab]{εγγραφή\\\code{authorized}} (gate.south |- k.north);
  \draw[->] (k.north -| an.south) -- node[right, lab]{\en{uevents}, ανάγνωση\\ περιγραφέων} (an.south);

  \node[draw, dashed, rounded corners, fit=(gate), inner sep=6pt, label={[lab]above:έμπιστη βάση (\en{TCB})}] {};
  \node[draw, dashed, rounded corners, fit=(an)(ag), inner sep=6pt, label={[lab]above:μη προνομιούχα τμήματα}] {};
\end{tikzpicture}
\caption{Αρχιτεκτονική του \Probolos{} με διαχωρισμό προνομίων. Μόνο η πύλη
εκτελείται ως \en{root}· ο αναλυτής λαμβάνει ανοιχτούς περιγραφείς αρχείων και
δεν έχει δικαίωμα εγγραφής στο \en{sysfs}.}
\label{fig:architecture}
\end{figure}

\section{Η διοχέτευση τεσσάρων σταδίων}
\label{sec:pipeline}

Κάθε νέα συσκευή εξετάζεται σε έως τέσσερα στάδια (Πίνακας~\ref{tab:stages},
Σχήμα~\ref{fig:pipeline}). Τα δύο πρώτα εκτελούνται ενώ η συσκευή είναι
πλήρως απενεργοποιημένη· τα στάδια 3 και 4 την ενεργοποιούν σύντομα και
ελεγχόμενα.

\begin{table}[htbp]
\centering
\caption{Τα τέσσερα στάδια εξέτασης και η κατάσταση της συσκευής σε καθένα.}
\label{tab:stages}
\begin{tabular}{@{}clp{4.2cm}p{4.4cm}@{}}
\toprule
\textbf{\#} & \textbf{Στάδιο} & \textbf{Ερώτημα} & \textbf{Κατάσταση συσκευής}\\
\midrule
1 & Ταυτότητα  & Τι ισχυρίζεται ότι είναι; & απενεργοποιημένη\\
2 & Συνέπεια   & Συμφωνούν οι ισχυρισμοί της; & απενεργοποιημένη\\
3 & Συμπεριφορά & Πληκτρολογεί μόνη της; & ενεργή για λίγα δευτερόλεπτα, είσοδος δεσμευμένη\\
4 & Αποθήκευση & Τι διαμερίσματα φέρει; & ενεργή σύντομα, μόνο ανάγνωση, χωρίς προσάρτηση\\
\bottomrule
\end{tabular}
\end{table}

\begin{figure}[htbp]
\centering
\begin{tikzpicture}[
  node distance=0.55cm,
  st/.style={draw, rounded corners=2pt, minimum width=2.25cm, minimum height=1.05cm,
             align=center, font=\footnotesize},
  dec/.style={draw, diamond, aspect=1.9, inner sep=1pt, align=center, font=\scriptsize, fill=gray!8},
  >={Stealth[length=2mm]}]
  \node[st, fill=gray!10] (plug) {Σύνδεση\\\code{authorized=0}};
  \node[st, fill=univblue!10, right=of plug] (s1) {Στάδιο 1\\Ταυτότητα};
  \node[st, fill=univblue!10, right=of s1] (s2) {Στάδιο 2\\Συνέπεια};
  \node[st, fill=warning!12, right=of s2] (s3) {Στάδιο 3\\Συμπεριφορά};
  \node[st, fill=warning!12, right=of s3] (s4) {Στάδιο 4\\Αποθήκευση};
  \node[dec, below=1.1cm of s3] (q) {Απόφαση\\χρήστη};
  \node[st, fill=okgreen!15, below left=0.9cm and 0.6cm of q] (ok) {\code{authorized=1}\\(μία φορά / πάντα)};
  \node[st, fill=critical!12, below right=0.9cm and 0.6cm of q] (no) {Παραμένει\\αποκλεισμένη};
  \node[st, fill=gray!8, below=1.1cm of s1, text width=2.3cm] (hist) {Ιστορικό \&\\εμπιστοσύνη};
  \draw[->] (plug) -- (s1);
  \draw[->] (s1) -- (s2);
  \draw[->] (s2) -- (s3);
  \draw[->] (s3) -- (s4);
  \draw[->] (s4) |- (q);
  \draw[->] (q) -| (ok);
  \draw[->] (q) -| (no);
  \draw[<->, dashed] (s2) |- (hist);
  \draw[->, dashed] (hist) |- node[pos=0.25, left, font=\scriptsize, align=center]{γνωστή \&\\ χωρίς \sevC} (ok);
\end{tikzpicture}
\caption{Ροή εξέτασης μιας νέας συσκευής. Μια απομνημονευμένη συσκευή
γίνεται δεκτή χωρίς ερώτηση, εκτός αν προκύψει εύρημα \sevC.}
\label{fig:pipeline}
\end{figure}

\subsection{Στάδιο 1: Ταυτότητα}
\label{sec:design-stage1}

Το στάδιο 1 διαβάζει ό,τι δηλώνει η συσκευή: τον περιγραφέα συσκευής, όλες τις
διαμορφώσεις, τις διεπαφές και τις συμβολοσειρές, την ταχύτητα σύνδεσης και την
κατανάλωση ισχύος. Η ανάγνωση γίνεται από το αρχείο \code{descriptors} του
\en{sysfs}, το οποίο περιέχει τα ακατέργαστα \en{bytes} που διάβασε ο πυρήνας
κατά την απαρίθμηση. Η συσκευή \emph{δεν} χρειάζεται να είναι εξουσιοδοτημένη.

Οι τύποι λειτουργιών μιας συσκευής προκύπτουν \emph{αποκλειστικά} από τις
κλάσεις των διεπαφών της και ποτέ από το όνομά της: μια συσκευή που ονομάζεται
«\en{SanDisk Cruzer}» αλλά εκθέτει διεπαφή \en{HID} πληκτρολογίου θεωρείται
πληκτρολόγιο.

\subsection{Στάδιο 2: Κανόνες συνέπειας}
\label{sec:design-stage2}

Το στάδιο 2 είναι η καρδιά της ανίχνευσης. Εφαρμόζει 21 κανόνες που εξετάζουν
αν οι ισχυρισμοί της συσκευής συνάδουν μεταξύ τους και με την προδιαγραφή. Οι
κανόνες ομαδοποιούνται σε τέσσερις οικογένειες:

\begin{description}
  \item[Αντιφατικοί συνδυασμοί λειτουργιών.] π.χ. \code{storage-with-keyboard}
        (\sevC): συσκευή αποθήκευσης που δηλώνει και πληκτρολόγιο·
        \code{network-with-keyboard} (\sevC)· \code{keyboard-with-unrelated-function}
        (\sevW).
  \item[Δομική ακεραιότητα περιγραφέων.] π.χ. \code{descriptor-chain-truncated}
        (\sevC): η αλυσίδα περιγραφέων είναι μικρότερη από ό,τι δηλώνει το
        \code{wTotalLength}· \code{configurations-missing} (\sevC)·
        \code{interface-count-mismatch} (\sevW).
  \item[Εχθρικό κείμενο.] π.χ. \code{crafted-strings-hid} (\sevC): συσκευή που
        μπορεί να πληκτρολογεί και κρύβει χαρακτήρες ελέγχου ή αμφίδρομης
        κατεύθυνσης (\en{bidi}) στο όνομά της~\cite{boucher2023trojansource,unicode-tr9}.
  \item[Φυσική αληθοφάνεια.] π.χ. \code{power-exceeds-bus-limit} (\sevW): ζητά
        περισσότερο ρεύμα από όσο επιτρέπει το \en{USB}·
        \code{storage-declares-negligible-power} (\sevN).
\end{description}

Ο πλήρης κατάλογος βρίσκεται στο Παράρτημα~\ref{app:rules}. Κάθε κανόνας που
καταρρέει κατά την εκτέλεση \emph{αναφέρεται} ως εύρημα
\code{analyzer-failed:<check>} και δεν παραλείπεται σιωπηλά: είναι \sevC{} όταν
ο έλεγχος καθορίζει το στάδιο του, \sevN{} αλλιώς.

\subsection{Στάδιο 3: Καραντίνα συμπεριφοράς}
\label{sec:design-stage3}

Η ιδέα της καραντίνας είναι απλή: η συσκευή ενεργοποιείται για λίγα δευτερόλεπτα
(προεπιλογή 3\,\en{s}), ενώ κάθε κόμβος εισόδου της δεσμεύεται αποκλειστικά με
\code{EVIOCGRAB}, και ο χρήστης έχει ενημερωθεί να \emph{μην την αγγίξει}. Κάθε
πληκτρολόγηση που φτάνει σε αυτό το διάστημα είναι εξ ορισμού ύποπτη.

Η σειρά των ενεργειών είναι κρίσιμη (Σχήμα~\ref{fig:quarantine}): η συσκευή
\emph{απενεργοποιείται ξανά πριν} απελευθερωθεί η δέσμευση, ακόμη και σε
περίπτωση εξαίρεσης. Αν αντίθετα απελευθερωνόταν πρώτα η δέσμευση, οι
πληκτρολογήσεις που βρίσκονταν στον ενδιάμεσο χώρο θα έφταναν στη συνεδρία.

\begin{figure}[htbp]
\centering
\begin{tikzpicture}[>={Stealth[length=2mm]}, font=\footnotesize]
  \draw[->, thick] (0,0) -- (13,0) node[right]{$t$};
  \foreach \x/\l in {0.6/{\code{authorized=1}}, 2.6/{πρώτο\\\code{EVIOCGRAB}}, 9.0/{\code{authorized=0}}, 11.8/{απελευθέρωση\\δέσμευσης}} {
    \draw (\x,0.15) -- (\x,-0.15);
    \node[below, align=center, text width=2.2cm] at (\x,-0.2) {\l};
  }
  \fill[critical!25] (0.6,0.25) rectangle (2.6,0.75);
  \node at (1.6,0.5) {κενό έκθεσης};
  \fill[okgreen!20] (2.6,0.25) rectangle (9.0,0.75);
  \node at (5.8,0.5) {παρατήρηση (\code{--observe}, προεπ. 3\,\en{s})};
  \fill[gray!20] (9.0,0.25) rectangle (11.8,0.75);
  \node at (10.4,0.5) {ασφαλές};
  \draw[decorate, decoration={brace, amplitude=4pt}] (0.6,0.85) -- (2.6,0.85)
     node[midway, above=4pt, align=center]{41–85\,\en{ms} σε πραγματικό υλικό};
\end{tikzpicture}
\caption{Χρονοδιάγραμμα της καραντίνας. Πληκτρολογήσεις που στέλνονται στο
«κενό έκθεσης», πριν από τη δέσμευση, μπορεί να διαφύγουν προς τη συνεδρία.}
\label{fig:quarantine}
\end{figure}

\paragraph{Κανόνες συμπεριφοράς.}
Έστω $n$ ο αριθμός των πατημάτων πλήκτρων και $\Delta_i = t_{i+1}-t_i$ τα
διαστήματα μεταξύ διαδοχικών πατημάτων, με χρονοσφραγίδες \emph{του πυρήνα}. Ο
μέσος χρόνος $\bar\Delta$ και ο συντελεστής μεταβλητότητας
\begin{equation}
  CV = \frac{\sigma_\Delta}{\bar\Delta},\qquad
  \sigma_\Delta = \sqrt{\frac{1}{n-1}\sum_{i=1}^{n-1}\left(\Delta_i-\bar\Delta\right)^2}
  \label{eq:cv}
\end{equation}
χρησιμοποιούνται για να διακριθεί η μηχανική από την ανθρώπινη πληκτρολόγηση.
Η συσκευή χαρακτηρίζεται \code{machine-generated-keystrokes} (\sevC) όταν
\begin{equation}
  n \geq 3 \;\wedge\; \bar\Delta < 50\,\text{\en{ms}} \;\wedge\; CV < 0{,}20 .
  \label{eq:machine}
\end{equation}
Διαφορετικά, αν $n \ge 5$, ενεργοποιείται ο κανόνας \code{unprompted-typing}
(\sevC), ενώ λιγότερες πληκτρολογήσεις δίνουν \code{unexpected-keystrokes}
(\sevW) — συμβατό τόσο με τυχαίο άγγιγμα όσο και με επίθεση. Τέλος, αν η πρώτη
πληκτρολόγηση φτάσει σε λιγότερο από 0{,}5\,\en{s} και $n\ge3$, προστίθεται
\code{immediate-activity} (\sevW).

\begin{remark}
Δεδομένου ότι ο ρυθμός μπορεί να πλαστογραφηθεί~\cite{farhi2019malboard}, ο
κανόνας \eqref{eq:machine} θεωρείται παρακάμψιμος. Ο κανόνας που \emph{μετρά}
είναι ο \code{unprompted-typing}: σε καραντίνα κανείς δεν πρέπει να πληκτρολογεί,
όσο ανθρώπινος κι αν μοιάζει ο ρυθμός.
\end{remark}

\subsection{Στάδιο 4: Επιθεώρηση αποθήκευσης}
\label{sec:design-stage4}

Για συσκευές \emph{αποκλειστικά} αποθήκευσης, το \Probolos{} διαβάζει τους
πρώτους τομείς του δίσκου ακατέργαστα και μόνο για ανάγνωση: τον πίνακα
\en{MBR} ή \en{GPT} και υπογραφές συστημάτων αρχείων (\en{NTFS}, \en{exFAT},
\en{FAT}, \en{ext2/3/4}, \en{btrfs}, \en{ISO~9660}, \en{UDF}). Ο δίσκος
\emph{δεν προσαρτάται ποτέ}. Η ανάγνωση εκτελείται σε ξεχωριστή διεργασία με
αυστηρό χρονικό όριο, ώστε μια συσκευή που «παγώνει» να μην παγώσει το
σύστημα.

Το στάδιο 4 παραλείπεται για σύνθετες συσκευές, διότι η ενεργοποίησή τους θα
ενεργοποιούσε και τις άλλες λειτουργίες τους. Το αποτέλεσμά του θεωρείται
\emph{πλαίσιο}: οι αποφάσεις στηρίζονται στα στάδια 1 και 2, αφού μια συσκευή
μπορεί να καθυστερήσει την εμφάνιση του μέσου της μέχρι να εγκαταλείψει ο έλεγχος.

\section{Μοντέλο σοβαρότητας και κλιμακούμενη τριβή}
\label{sec:design-friction}

Τα ευρήματα έχουν τρία επίπεδα σοβαρότητας:

\begin{description}
  \item[\sevC] Η συσκευή ταιριάζει με πρότυπο επίθεσης. Η αποδοχή απαιτεί
        αντίστροφη μέτρηση 10 δευτερολέπτων στον διάλογο ή πληκτρολόγηση της λέξης
        \code{authorize} στο τερματικό· η επιλογή «πάντα» δεν προσφέρεται και η
        αποθηκευμένη εμπιστοσύνη δεν την παρακάμπτει.
  \item[\sevW] Ύποπτη συσκευή.
  \item[\sevN] Αξιοσημείωτη πληροφορία.
\end{description}

Η αλληλεπίδραση με τον χρήστη ακολουθεί την αρχή της ψυχολογικής
αποδοχής~\cite{saltzer1975protection}: η τριβή είναι ανάλογη του κινδύνου
(Πίνακας~\ref{tab:friction}). Μια συνηθισμένη συσκευή αποθήκευσης χωρίς ευρήματα
χρειάζεται ένα κλικ· μια συσκευή που μπορεί να πληκτρολογεί χρειάζεται δύο· μια
συσκευή με κρίσιμο εύρημα απαιτεί αναμονή. Επιπλέον, η αντίστροφη μέτρηση
\emph{επιβάλλεται από τον δαίμονα}: μια έγκριση που φτάνει νωρίτερα από 10
δευτερόλεπτα απορρίπτεται, ανεξάρτητα από τον αποστολέα.

\begin{table}[htbp]
\centering
\caption{Κλιμακούμενη τριβή των διαλόγων επιφάνειας εργασίας.}
\label{tab:friction}
\begin{tabular}{@{}p{6.4cm}p{7.2cm}@{}}
\toprule
\textbf{Συσκευή} & \textbf{Αλληλεπίδραση}\\
\midrule
Δεν μπορεί να πληκτρολογεί ή να μεταφέρει κίνηση, χωρίς προειδοποιήσεις & \textbf{ένας} διάλογος\\
\en{HID}, δίκτυο, ραδιοσυχνότητες, άγνωστη ή \en{vendor-specific}, ή οποιαδήποτε \sevW & \textbf{δύο} διάλογοι\\
Οποιοδήποτε εύρημα \sevC & \textbf{αντίστροφη μέτρηση} 10\,\en{s} πριν ενεργοποιηθεί το «Αποδοχή παρ' όλα αυτά»\\
\bottomrule
\end{tabular}
\end{table}

\paragraph{Η συσκευή δεν μπορεί να απαντήσει για τον εαυτό της.}
Επειδή η ερώτηση εμφανίζεται ενώ η συσκευή είναι απενεργοποιημένη, ένα ψεύτικο
πληκτρολόγιο δεν μπορεί να πατήσει «Ναι». Επιπλέον, η ερώτηση στο τερματικό
\emph{αδειάζει πρώτα} την εκκρεμή είσοδο και αναφέρει πόσα \en{bytes} απέρριψε,
ώστε πλήκτρα που μπορεί να διέφυγαν κατά το στάδιο 3 να μην απαντήσουν στη θέση
του χρήστη.

\section{Ιστορικό, εμπιστοσύνη και μετάλλαξη}
\label{sec:design-trust}

Το σύστημα διατηρεί δύο ξεχωριστές αποθήκες:

\begin{itemize}
  \item \textbf{Βιβλίο ιστορικού} (\en{ledger}): κάθε συσκευή που έχει δει, ο
        τύπος της, οι αποφάσεις και ένα \emph{κανονικοποιημένο} αποτύπωμα.
  \item \textbf{Χώρος εμπιστοσύνης} (\en{trust store}): συσκευές που εγκρίθηκαν
        «πάντα», δεσμευμένες με την ταυτότητα \emph{και} τα ακριβή
        \en{bytes} των περιγραφέων.
\end{itemize}

\begin{definition}[Κανονικοποιημένο αποτύπωμα]
Έστω $D$ οι αναλυμένοι περιγραφείς μιας συσκευής. Το αποτύπωμα του ιστορικού
είναι
\[
  F(D) = \mathrm{SHA256}\big(\mathit{VID}\,\|\,\mathit{PID}\,\|\,\mathit{bcdDevice}\,\|\,
         \mathit{class}_{dev}\,\|\,\textstyle\prod_{c,i}(n_c,\mathit{class}_{c,i})\big),
\]
όπου το γινόμενο συμβολίζει τη διαδοχική παράθεση, για κάθε διαμόρφωση $c$ και
διεπαφή $i$, των πλήθων και της τριάδας κλάσης/υποκλάσης/πρωτοκόλλου.
Παραλείπονται σκόπιμα τα \code{bcdUSB}, \code{bMaxPacketSize0},
\code{bMaxPower}, οι περιγραφείς τελικών σημείων και οι συνοδευτικοί περιγραφείς
\en{SuperSpeed}, τα οποία αλλάζουν νόμιμα όταν μια συσκευή μετακινείται μεταξύ
θύρας \en{USB}~2 και \en{USB}~3.
\end{definition}

Ο κανόνας \code{descriptor-drift} (\sevC) ενεργοποιείται όταν μια συσκευή με
γνωστή ταυτότητα (\en{VID}, \en{PID}, σειριακός αριθμός) παρουσιάζει διαφορετικό
$F(D)$. Αυτό ανιχνεύει το σενάριο όπου ένα «αθώο» στικάκι εγκρίνεται και αργότερα
επαναπρογραμματίζεται ώστε να εκθέτει και πληκτρολόγιο. Ο κανόνας
\code{previously-rejected} (\sevC) ενεργοποιείται για συσκευή που ο χρήστης
απέρριψε στο παρελθόν.

Ο χώρος εμπιστοσύνης είναι αυστηρότερος: το κλειδί του είναι η ταυτότητα και ένα
\en{SHA-256} των \emph{ακατέργαστων} \en{bytes}. Έτσι μια απομνημονευμένη
συσκευή που μετακινείται σε άλλο τύπο θύρας ερωτάται ξανά — αποτυχία προς την
ασφαλή κατεύθυνση. Ένας χώρος εμπιστοσύνης που δεν μπορεί να διαβαστεί ή
περιέχει έστω και μία κακοσχηματισμένη εγγραφή \emph{δεν εμπιστεύεται τίποτα}.

\section{Διαχωρισμός προνομίων}
\label{sec:design-privsep}

Ακολουθώντας την αρχή του ελάχιστου προνομίου~\cite{provos2003privsep}, ο
\code{--privsep} χωρίζει το \Probolos{} σε δύο διεργασίες που επικοινωνούν μέσω
ιδιωτικού \code{SEQPACKET} \en{socketpair}, δημιουργημένου πριν από το
\code{fork()}. Η πύλη περνά ανοιχτούς περιγραφείς αρχείων με \code{SCM\_RIGHTS}.
Οι βασικές αρχές είναι:

\begin{enumerate}
  \item \textbf{Το πεδίο δράσης προκύπτει από τον πυρήνα, όχι από το αίτημα.} Η
        πύλη ενεργεί μόνο σε συσκευή \en{USB} με \code{authorized=0}, ή σε
        συσκευή που η ίδια ενεργοποίησε για επιθεώρηση, και μόνο στη
        συγκεκριμένη \emph{στιγμιότυπη εμφάνιση} (\en{instance}) του πυρήνα. Μια
        συσκευή που επανα-απαριθμήθηκε στην ίδια θύρα είναι άλλο στιγμιότυπο
        και απορρίπτεται.
  \item \textbf{Επαληθευμένη απόρριψη προνομίων.} Αφαιρούνται οι συμπληρωματικές
        ομάδες (μια ομάδα \code{input} ή \code{disk} που θα επιβίωνε θα
        παρέκαμπτε εντελώς την πύλη), έπειτα \en{gid} και \en{uid}, και ελέγχεται
        ότι το \code{setuid(0)} \emph{αποτυγχάνει}. Ο αναλυτής καλεί
        \code{setsid()} ώστε να μην έχει τερματικό ελέγχου στο οποίο θα
        μπορούσε να εγχύσει πλήκτρα (\code{TIOCSTI}).
  \item \textbf{Η «μόνιμη εμπιστοσύνη» γράφεται μόνο από την πύλη}, μόνο για
        συσκευή που η ίδια ενεργοποίησε τα τελευταία 60\,\en{s}, και μόνο με
        αποτύπωμα που \emph{μέτρησε η ίδια} πριν από την πρώτη εγγραφή
        \code{authorized=1}. Ο αναλυτής μπορεί να ζητήσει, αλλά το κλειδί του
        πρέπει να ταυτίζεται ακριβώς.
  \item \textbf{Τίποτα δεν επιβιώνει του αναλυτή.} Η διεργασία \en{root} είναι
        \en{child subreaper} (\code{PR\_SET\_CHILD\_SUBREAPER}): κάθε διεργασία που
        αφήνει πίσω του ο αναλυτής τερματίζεται.
\end{enumerate}

Το Σχήμα~\ref{fig:trust-seq} δείχνει τη σειρά μηνυμάτων για μια έγκριση «πάντα».

\begin{figure}[htbp]
\centering
\begin{tikzpicture}[font=\footnotesize, >={Stealth[length=2mm]},
  note/.style={draw, fill=#1, rounded corners=2pt, align=center, font=\scriptsize, inner sep=3pt},
  note/.default=gray!8]
  \foreach \x/\n in {0/Πυρήνας, 4/Πύλη (\en{root}), 8.4/Αναλυτής, 12.4/Πράκτορας} {
    \node[draw, fill=gray!10, minimum width=2.4cm, align=center] at (\x,0) {\n};
    \draw[dashed, gray] (\x,-0.35) -- (\x,-9.0);
  }
  \draw[->] (0,-0.9) -- node[above, pos=0.3]{\en{uevent} \code{add}} (8.4,-0.9);
  \node[note, right] at (8.5,-1.55) {στάδια 1–2};
  \draw[->] (8.4,-2.2) -- node[above]{\code{ASK}} (12.4,-2.2);
  \draw[->] (12.4,-2.9) -- node[above]{«Αποδοχή, πάντα»} (8.4,-2.9);
  \draw[->] (8.4,-3.6) -- node[above]{\code{AUTHORIZE(instance)}} (4,-3.6);
  \node[note=critical!8, left] at (3.9,-4.15) {μέτρηση\\αποτυπώματος $K_g$};
  \draw[->] (4,-5.0) -- node[above, pos=0.6]{\code{authorized=1}} (0,-5.0);
  \draw[->] (8.4,-5.8) -- node[above]{\code{TRUST(instance}, $K_a$\code{)}} (4,-5.7);
  \node[note=okgreen!12, right] at (4.1,-6.7) {έλεγχος: $\le 60$\,\en{s}, ίδιο στιγμιότυπο,\\ $K_a = K_g$, μία χρήση};
  \draw[->] (4,-7.8) -- node[above]{εγγραφή \code{trusted.json}} (0,-7.8);
  \draw[->] (4,-8.5) -- node[above]{\code{OK}} (8.4,-8.5);
\end{tikzpicture}
\caption{Ακολουθία μηνυμάτων για έγκριση «πάντα» υπό \code{--privsep}. Ο
αναλυτής δεν μπορεί να επιβάλει αποτύπωμα· η πύλη συγκρίνει με το δικό της.}
\label{fig:trust-seq}
\end{figure}

\section{Ασφάλεια έναντι κλειδώματος}
\label{sec:design-failsafe}

Ο μεγαλύτερος κίνδυνος ενός εργαλείου προεπιλεγμένης άρνησης είναι να αφήσει
τον χρήστη χωρίς πληκτρολόγιο. Η επαναφορά του \code{authorized\_default}
συνεπώς δεν είναι απλώς ρουτίνα τερματισμού, αλλά συνδέεται σε τέσσερις
ανεξάρτητες διαδρομές: (1) τον διαχειριστή περιβάλλοντος (\en{context manager})
της \en{Python}, (2) χειριστές σημάτων \code{SIGTERM}, \code{SIGINT},
\code{SIGHUP}, (3) το \code{atexit}, και (4) το μήνυμα ανάκτησης. Μόνο το
\code{SIGKILL} και ένας \en{kernel panic} μπορούν να τις παρακάμψουν. Επιπλέον:

\begin{itemize}
  \item \textbf{Πρώτα εγγραφές, μετά λόγια.} Κάθε διανομέας επαναφέρεται
        \emph{πριν} καταγραφεί οποιοδήποτε μήνυμα· ένα μήνυμα που δεν μπορεί να
        τυπωθεί (π.χ. κλειστό τερματικό) δεν σταματά τίποτα.
  \item \textbf{\en{Watchdog}:} αν ο δαίμονας σταματήσει να προοδεύει για
        \code{--watchdog} δευτερόλεπτα (προεπ. 60), η πύλη ανοίγει και η
        διεργασία τερματίζεται με κωδικό αποτυχίας ώστε το \en{systemd} να την
        επανεκκινήσει.
  \item \textbf{Αρχείο πανικού:} \code{sudo touch /run/probolos.panic} ανοίγει την
        πύλη· το αρχείο πρέπει να ανήκει στον \en{root}, ώστε κανένας χρήστης να
        μην μπορεί να απενεργοποιήσει την προστασία.
  \item \textbf{Εξαιρέσεις θυρών:} οι εσωτερικές θύρες (\code{removable=fixed})
        και όσες δηλωθούν με \code{--allow-port} δεν ελέγχονται, ούτε όταν η
        οθόνη είναι κλειδωμένη, ώστε ένα πληκτρολόγιο διάσωσης να μπορεί να
        πληκτρολογήσει τον κωδικό.
  \item \textbf{Υπολείμματα:} αν κατά την εκκίνηση βρεθεί \code{authorized\_default=0}
        (υπόλειμμα κατάρρευσης), δεν θεωρείται αρχική τιμή· επαναφέρεται το 1.
\end{itemize}

\section{Εχθρικό κείμενο συσκευών}
\label{sec:design-text}

Μια συσκευή επιλέγει ελεύθερα τις συμβολοσειρές κατασκευαστή, προϊόντος και
σειριακού αριθμού, οι οποίες φτάνουν στο τερματικό, στους διαλόγους, στο αρχείο
καταγραφής \en{JSON} και στις αποθήκες. Για να μην μπορεί μια συσκευή να
«ξαναγράψει» την αναφορά για τον εαυτό της (π.χ. με ακολουθίες διαφυγής
\en{ANSI} ή χαρακτήρες \en{bidi}~\cite{boucher2023trojansource}), το κείμενο
καθαρίζεται σε \emph{ένα μόνο σημείο}: εκεί όπου τα \en{bytes} του \en{sysfs}
γίνονται συμβολοσειρές \en{Python}. Χαρακτήρες ελέγχου, \en{bidi} και αόρατοι
χαρακτήρες εμφανίζονται με διαφυγή και ο \emph{λόγος} του καθαρισμού γίνεται
εύρημα. Για τους διαλόγους \en{kdialog} (\en{Qt rich text}) και \en{zenity}
(\en{Pango}) εφαρμόζεται επιπλέον διαφυγή σήμανσης.
